\scanpage{92}
l'homomorphisme canonique
\[
\tag{45}\mathcal C'\longrightarrow\widetilde{\mathcal C}^{\unsure{\prime|*}}.
\]
\inserted{vrai revêtement de $\mathcal C$, dont les sommets sont les $\pi_0$ des fibres de $\mathcal C'-\unsure{\mathcal C|\widetilde{\mathcal C}}$}
\deleted{\unsure{ce}} foncteur $\mathcal C'-\mathcal C$ est injectif, de sorte qu'on peut identifier $\mathcal C'$ à un sous-groupoïde plein de $\widetilde{\mathcal C}^{\unsure{\prime|*}}$, équivalent à $\widetilde{\mathcal C}^{\unsure{\prime|*}}$ par l'inclusion, où $\widetilde{\mathcal C}$ est un rev. fini de $\mathcal C$.

Soit $x'_0$ un sommet de $\mathcal C'$, $x_0$ son image dans $\widetilde{\mathcal C}$. Soit $\Theta$ un automorphisme de $\mathcal C$ qui fixe $x_0$, \deleted{alors s'}il \deleted{\unsure{y en a un de}} \inserted{revient au} $=$ de se donner un autom. de \deleted{$\unsure{\mathcal C}$} $\pi_1(\mathcal C,x_0)\overset{\mathrm{def}}{=}\pi_0$ (\unsure{abrég.}), et pour tout \inserted{sommet} $\deleted{\unsure{x}}x_i\in\operatorname{Ob}\mathcal C$, $x_i\ne x_0$, un automorphisme du $\pi_1(x_0)$-torseur à droite $H_{x_i,x_0}=\operatorname{Hom}_{\mathcal C}(x_0,x_i)$, compatible avec l'opération de $\Theta$ sur $\pi_0$. Ceci posé, pour que cette action se remonte en une action de $\widetilde{\mathcal C}$ qui fixe $x'_0$, il f. et s. que le sous-groupe $\pi'_0$ de $\pi_0$, image de \deleted{$\unsure{\pi_1}$}
\[
\pi_1(\widetilde{\mathcal C},x'_0)\hookrightarrow\pi_1(\mathcal C,x_0)=\pi_0,
\]
soit \uline{stable} par $\Theta$ --- l'action de $\Theta$ sur les fibres \inserted{$\widetilde{\mathcal C}_{x_i}$} de $\widetilde{\mathcal C}/\mathcal C$ est alors donnée par la formule
\[
\tag{45}\operatorname{Ob}\widetilde{\mathcal C}_{x_i}=H_{x_i,x_0}\wedge_{\pi_0}\underbrace{\pi_0/\pi'_0}_{\substack{\text{esp. hom}\\\text{des classes}\\\text{à droite, sur lequel}\\\pi_0\ \text{opère à g}}}
\]
par transport de structure à partir de son
\margin{91}
